% ================================================================
%  APPENDICES
%  Status: DRAFT. Appendix B populated from sweeps2/ (tau-uniform
%  sweep) + floors_and_tables.py (Gibbs floors).
% ================================================================

\chapter{Derivations}
\label{app:derivations}

% ----------------------------------------------------------------
\section{Detailed balance for Metropolis--Hastings}
\label{app:mh-detailed-balance}

This appendix verifies the claim of Section~\ref{sec:metropolis-hastings}
that the Metropolis--Hastings chain satisfies detailed balance with
respect to the posterior $\pi(\bm{\theta}) = p(\bm{\theta} \given
\mathbf{y})$. The transition kernel has an accepted-move component and
a rejected-move component,
%
\begin{equation}
  T(\bm{\theta}' \given \bm{\theta})
  = q(\bm{\theta}' \given \bm{\theta})\,
    \alpha(\bm{\theta}, \bm{\theta}')
  + \delta(\bm{\theta}' - \bm{\theta})\,r(\bm{\theta}),
\end{equation}
%
where $\delta$ is the Dirac delta and
$r(\bm{\theta}) = \int [1 - \alpha(\bm{\theta},
\bm{\theta}'')]\,q(\bm{\theta}'' \given \bm{\theta})
\dd\bm{\theta}''$ is the total rejection probability. For
$\bm{\theta}' = \bm{\theta}$ detailed balance holds trivially, so take
$\bm{\theta}' \neq \bm{\theta}$, where the delta term vanishes.
Substituting the acceptance probability~\eqref{eq:mh-acceptance},
%
\begin{align}
  T(\bm{\theta}' \given \bm{\theta})\,\pi(\bm{\theta})
  &= q(\bm{\theta}' \given \bm{\theta})\,
     \pi(\bm{\theta})\,
     \min\!\left(1,\;
     \frac{\pi(\bm{\theta}')\,q(\bm{\theta} \given \bm{\theta}')}
          {\pi(\bm{\theta})\,q(\bm{\theta}' \given \bm{\theta})}
     \right)
     \notag \\[2pt]
  &= \min\!\Big(
     q(\bm{\theta}' \given \bm{\theta})\,\pi(\bm{\theta}),\;
     q(\bm{\theta} \given \bm{\theta}')\,\pi(\bm{\theta}')
     \Big)
     \notag \\[2pt]
  &= q(\bm{\theta} \given \bm{\theta}')\,\pi(\bm{\theta}')\,
     \min\!\left(
     \frac{\pi(\bm{\theta})\,q(\bm{\theta}' \given \bm{\theta})}
          {\pi(\bm{\theta}')\,q(\bm{\theta} \given \bm{\theta}')},\;
     1\right)
     \notag \\[2pt]
  &= T(\bm{\theta} \given \bm{\theta}')\,\pi(\bm{\theta}'),
\end{align}
%
where the second equality moves the positive factor
$q(\bm{\theta}' \given \bm{\theta})\,\pi(\bm{\theta})$ inside the
minimum and the third moves the positive factor
$q(\bm{\theta} \given \bm{\theta}')\,\pi(\bm{\theta}')$ back out. The
second line is symmetric under exchange of $\bm{\theta}$ and
$\bm{\theta}'$, which is what the third line exploits. Detailed
balance holds, so $\pi$ is the stationary distribution of the chain.

% ----------------------------------------------------------------
\section{The optimal importance sampling proposal}
\label{app:is-optimal-proposal}

This appendix proves the claim of Section~\ref{sec:is-basics} that, among
all proposals $q$, the variance of the importance sampling estimator
$\hat{L}_{\mathrm{IS}}
= \frac{1}{N}\sum_{i} w(\bm{\theta}_i)\,f(\bm{\theta}_i)$, with weight
$w(\bm{\theta}) = p(\bm{\theta})/q(\bm{\theta})$, is minimised by
%
\begin{equation}
  q^*(\bm{\theta}) \propto |f(\bm{\theta})|\, p(\bm{\theta}).
  \label{eq:app-is-optimal}
\end{equation}

\begin{proof}
Since $\hat{L}_{\mathrm{IS}}$ is unbiased for every $q$, minimising its
variance is equivalent to minimising the second moment of a single term,
\begin{equation}
  \E_q\!\left[\left(\frac{f(\bm{\theta}) p(\bm{\theta})}{q(\bm{\theta})}\right)^2\right]
  = \int \frac{f(\bm{\theta})^2 p(\bm{\theta})^2}{q(\bm{\theta})}\, \dd\bm{\theta}.
\end{equation}
By the Cauchy--Schwarz inequality, for any density $q$,
\begin{equation}
  \int \frac{f^2 p^2}{q}\,\dd\bm{\theta}
  = \left(\int \frac{f^2 p^2}{q}\,\dd\bm{\theta}\right)\!\!\left(\int q\,\dd\bm{\theta}\right)
  \geq \left(\int |f|\, p\, \dd\bm{\theta}\right)^2,
\end{equation}
with equality if and only if $f^2 p^2 / q \propto q$, that is,
\begin{equation}
  q^*(\bm{\theta}) \propto |f(\bm{\theta})|\, p(\bm{\theta}).
\end{equation}
The lower bound $\left(\int |f| p\, \dd\bm{\theta}\right)^2$ is independent of $q$,
so $q^*$ attains it and is therefore optimal. When $f \geq 0$ the resulting
weighted integrand $w(\bm{\theta}) f(\bm{\theta})$ is constant, and the variance is zero.
\end{proof}

% ================================================================
\chapter{Experimental details}
\label{app:experimental}

This appendix records the configuration and per-cell results of the
Bayesian linear regression experiment of
Chapter~\ref{ch:rao-blackwell}.

% ----------------------------------------------------------------
\section{Configuration and protocol}
\label{app:config}

The model is the hierarchical linear regression of
equation~\eqref{eq:lr-model}, with prior scale $\sigma_0 = 5$ and
noise-precision prior $\tau \sim \Uniform(0.01, 1)$, placing the noise
standard deviation in $\sigma \in (1, 10)$. The estimator is a
multilayer perceptron of depth $3$ and hidden width $128$ with ReLU
activations and a softplus output on the $\sigma$-head, trained with
Adam at learning rate $10^{-3}$.

The grid sweeps parameter dimension $p \in \{2, 5, 10, 20, 30\}$ and
minibatch size $B \in \{4, 16, 64, 256\}$ at sample size $n = 50$, with
a single design matrix fixed per dimension (a second, independent
realisation is examined in Appendix~\ref{app:second-design}). Each cell
is trained for $200$ epochs over ten seeds, the Monte Carlo and
Rao-Blackwellised estimators sharing their initialisation and their
simulated batches.
Each Bayes-optimal floor is computed by the Gibbs sampler of
Section~\ref{sec:bayes-floor} on a held-out test set of $1024$
problems; its variance-based and squared-error estimators agree at
every dimension (Table~\ref{tab:floors}).

% ----------------------------------------------------------------
\section{Main grid results}
\label{app:main-grid}

Table~\ref{tab:main-grid} reports, for every $(p, B)$ cell, the final
test loss of each estimator (the mean of the last ten epochs, averaged
over ten seeds), the Bayes floor $L^{*}$, the raw ratio
$\mathcal{L}_{\MC}/\mathcal{L}_{\RB}$, and the excess-over-floor ratio
$(\mathcal{L}_{\MC}-L^{*})/(\mathcal{L}_{\RB}-L^{*})$. The
Rao-Blackwellised estimator attains the lower loss in all twenty cells.

\begin{table}[htbp]
  \centering
  \small
  \caption{Main grid: final test loss of the Monte Carlo (MC) and
  Rao-Blackwellised (RB) estimators, the Bayes floor $L^{*}$, and the
  raw and excess-over-floor ratios, for each $(p, B)$ cell at $n = 50$.}
  \label{tab:main-grid}
  \begin{tabular}{rr rrr rr}
    \toprule
    $p$ & $B$ & $\mathcal{L}_{\MC}$ & $\mathcal{L}_{\RB}$
        & $L^{*}$ & ratio & excess ratio \\
    \midrule
    2 & 4   & 0.1080 & 0.0921 & 0.0708 & 1.172 & 1.74 \\
    2 & 16  & 0.1033 & 0.0906 & 0.0708 & 1.140 & 1.64 \\
    2 & 64  & 0.0979 & 0.0868 & 0.0708 & 1.129 & 1.70 \\
    2 & 256 & 0.0932 & 0.0864 & 0.0708 & 1.079 & 1.43 \\
    \midrule
    5 & 4   & 0.1549 & 0.1234 & 0.0964 & 1.255 & 2.16 \\
    5 & 16  & 0.1526 & 0.1190 & 0.0964 & 1.283 & 2.49 \\
    5 & 64  & 0.1410 & 0.1158 & 0.0964 & 1.218 & 2.30 \\
    5 & 256 & 0.1312 & 0.1225 & 0.0964 & 1.071 & 1.33 \\
    \midrule
    10 & 4   & 0.2209 & 0.1701 & 0.1102 & 1.299 & 1.85 \\
    10 & 16  & 0.1944 & 0.1484 & 0.1102 & 1.310 & 2.20 \\
    10 & 64  & 0.1700 & 0.1448 & 0.1102 & 1.174 & 1.73 \\
    10 & 256 & 0.1638 & 0.1622 & 0.1102 & 1.010 & 1.03 \\
    \midrule
    20 & 4   & 0.3766 & 0.2683 & 0.1263 & 1.404 & 1.76 \\
    20 & 16  & 0.2759 & 0.1983 & 0.1263 & 1.392 & 2.08 \\
    20 & 64  & 0.2104 & 0.1802 & 0.1263 & 1.168 & 1.56 \\
    20 & 256 & 0.2407 & 0.2328 & 0.1263 & 1.034 & 1.07 \\
    \midrule
    30 & 4   & 0.6366 & 0.5135 & 0.2188 & 1.240 & 1.42 \\
    30 & 16  & 0.4524 & 0.3362 & 0.2188 & 1.346 & 1.99 \\
    30 & 64  & 0.3806 & 0.3354 & 0.2188 & 1.135 & 1.39 \\
    30 & 256 & 0.3681 & 0.3569 & 0.2188 & 1.031 & 1.08 \\
    \bottomrule
  \end{tabular}
\end{table}

% ----------------------------------------------------------------
\section{Batch-size sweep}
\label{app:batch-sweep}

Table~\ref{tab:batch-sweep} reports the $p = 20$ cell at batch sizes
$B \in \{4, 16, 64, 256\}$ (ten seeds each, $200$ epochs). The loss
is the mean over the last ten epochs; the excess is that loss minus
the floor $L^{*} = 0.126$.

\begin{table}[htbp]
  \centering
  \small
  \caption{Batch-size sweep at $p = 20$: final loss and excess over the
  floor for each estimator.}
  \label{tab:batch-sweep}
  \begin{tabular}{r rr rr r}
    \toprule
    $B$ & MC loss & RB loss & MC excess & RB excess & excess ratio \\
    \midrule
    4   & 0.377 & 0.268 & 0.250 & 0.142 & 1.76 \\
    16  & 0.276 & 0.198 & 0.150 & 0.072 & 2.08 \\
    64  & 0.210 & 0.180 & 0.084 & 0.054 & 1.56 \\
    256 & 0.241 & 0.233 & 0.114 & 0.107 & 1.07 \\
    \bottomrule
  \end{tabular}
\end{table}

The Rao-Blackwellised estimator at batch $4$ (excess $0.142$) matches
the Monte Carlo estimator at batch $16$ (excess $0.150$): its variance
reduction is worth a fourfold increase in batch size. At every batch
size its excess over the floor is the smaller of the two, and only at
$B = 256$ (where both estimators have nearly exhausted the Monte
Carlo noise) do the two nearly coincide.

% ----------------------------------------------------------------
\section{Bayes-optimal floors}
\label{app:floors}

Table~\ref{tab:floors} reports the Bayes floor for each dimension at
$n = 50$, together with the two independent estimators of it: the
average within-chain posterior variance and the mean squared error of
the posterior-mean estimate (Section~\ref{sec:bayes-floor}). The two
agree to within sampling error at every dimension, validating the Gibbs
sampler.

\begin{table}[htbp]
  \centering
  \small
  \caption{Bayes-optimal floors by dimension $p$ (at $n = 50$), with
  the squared-error and variance-based estimators.}
  \label{tab:floors}
  \begin{tabular}{r rr}
    \toprule
    $p$ & $L^{*}$ (MSE) & $L^{*}$ (variance) \\
    \midrule
    2  & 0.0708 & 0.0774 \\
    5  & 0.0964 & 0.0976 \\
    10 & 0.1102 & 0.1131 \\
    20 & 0.1263 & 0.1276 \\
    30 & 0.2188 & 0.2165 \\
    \bottomrule
  \end{tabular}
\end{table}

% ----------------------------------------------------------------
\section{Second design matrix}
\label{app:second-design}

The main-text results fix a single design matrix $X$ per dimension. To
confirm that the Rao-Blackwellisation advantage is a property of the
estimator and not of one particular design, the entire sweep was
repeated on a second, independent design-matrix realisation $X^{(2)}$
(a different random seed for $X$; identical prior, grid, ten training
seeds, and $200$ epochs). Its Bayes floors are recomputed by the same
Gibbs sampler on $X^{(2)}$ and reported in the $L^{*}$ column of
Table~\ref{tab:main-grid-X2}.

The pattern of the main text reappears (Figure~\ref{fig:curves-X2}):
the Rao-Blackwellised estimator attains the lower loss in eighteen of
the twenty cells, by margins that grow as the batch size shrinks. The
two exceptions are both at the largest batch $B = 256$, where the two
estimators have all but exhausted the Monte Carlo noise and coincide to
within $2\%$ --- the same vanishing of the Rao-Blackwell advantage with
batch size seen for the first design. The advantage is therefore
robust to the design wherever it is not already exhausted by averaging.

\begin{figure}[htbp]
  \centering
  \IfFileExists{figures/rb_tau_curves_X2.png}
    {\includegraphics[width=\textwidth]{rb_tau_curves_X2}}
    {\fbox{\parbox[c][6cm][c]{\textwidth}{\centering\itshape
      Figure pending: second design-matrix curves grid
      (re-export \texttt{rb\_tau\_curves\_X2} via
      \texttt{replot\_curves\_floor.py}).}}}
  \caption{As Figure~\ref{fig:curves}, for the second design-matrix
  realisation $X^{(2)}$: training curves across parameter dimension $p$
  (rows) and batch size (columns); Monte Carlo (blue) and
  Rao-Blackwellised (orange), with the Bayes floor $L^{*}(p)$ dashed.}
  \label{fig:curves-X2}
\end{figure}

\begin{table}[htbp]
  \centering
  \small
  \caption{Second design matrix $X^{(2)}$: final test loss of the Monte
  Carlo (MC) and Rao-Blackwellised (RB) estimators, the Bayes floor
  $L^{*}$, and the raw and excess-over-floor ratios, for each $(p, B)$
  cell at $n = 50$.}
  \label{tab:main-grid-X2}
  \begin{tabular}{rr rrr rr}
    \toprule
    $p$ & $B$ & $\mathcal{L}_{\MC}$ & $\mathcal{L}_{\RB}$
        & $L^{*}$ & ratio & excess ratio \\
    \midrule
    2 & 4   & 0.1041 & 0.0882 & 0.0666 & 1.180 & 1.74 \\
    2 & 16  & 0.1001 & 0.0861 & 0.0666 & 1.162 & 1.72 \\
    2 & 64  & 0.0936 & 0.0838 & 0.0666 & 1.117 & 1.57 \\
    2 & 256 & 0.0889 & 0.0823 & 0.0666 & 1.080 & 1.42 \\
    \midrule
    5 & 4   & 0.1508 & 0.1234 & 0.0939 & 1.222 & 1.93 \\
    5 & 16  & 0.1471 & 0.1187 & 0.0939 & 1.240 & 2.15 \\
    5 & 64  & 0.1396 & 0.1160 & 0.0939 & 1.204 & 2.07 \\
    5 & 256 & 0.1271 & 0.1187 & 0.0939 & 1.070 & 1.34 \\
    \midrule
    10 & 4   & 0.2002 & 0.1561 & 0.0972 & 1.282 & 1.75 \\
    10 & 16  & 0.1762 & 0.1360 & 0.0972 & 1.295 & 2.04 \\
    10 & 64  & 0.1540 & 0.1294 & 0.0972 & 1.189 & 1.76 \\
    10 & 256 & 0.1524 & 0.1531 & 0.0972 & 0.995 & 0.99 \\
    \midrule
    20 & 4   & 0.3888 & 0.2864 & 0.1407 & 1.357 & 1.70 \\
    20 & 16  & 0.2912 & 0.2156 & 0.1407 & 1.351 & 2.01 \\
    20 & 64  & 0.2311 & 0.1953 & 0.1407 & 1.183 & 1.66 \\
    20 & 256 & 0.2538 & 0.2588 & 0.1407 & 0.981 & 0.96 \\
    \midrule
    30 & 4   & 0.7288 & 0.5779 & 0.2737 & 1.261 & 1.50 \\
    30 & 16  & 0.5409 & 0.3974 & 0.2737 & 1.361 & 2.16 \\
    30 & 64  & 0.4639 & 0.4018 & 0.2737 & 1.155 & 1.49 \\
    30 & 256 & 0.4435 & 0.4278 & 0.2737 & 1.037 & 1.10 \\
    \bottomrule
  \end{tabular}
\end{table}

% ----------------------------------------------------------------
\section{Reproduction}
\label{app:reproduction}

The experiment code is in \texttt{sweeps2/} of the accompanying
repository. The training curves of Figure~\ref{fig:curves} and the raw
per-cell trajectories are produced by
%
\begin{quote}
\texttt{python rb\_vs\_mc\_sweep\_tau.py}
\end{quote}
%
and the Bayes floors and the tables of this appendix by
\texttt{floors\_and\_tables.py}, which runs a vectorised Gibbs sampler
for the linear-regression posterior on the same held-out test sets. The
Bayes-floor line on Figure~\ref{fig:curves} is overlaid by
\texttt{replot\_curves\_floor.py}, and the samples-seen view of
Figure~\ref{fig:vs-samples} by \texttt{plot\_curves\_vs\_samples.py};
both reuse the saved \texttt{.npz} without retraining.
